\documentclass[12pt]{article}
\usepackage[utf8]{inputenc}
\usepackage[T1]{fontenc}
\usepackage[margin=1in]{geometry}
\usepackage{setspace}
\usepackage{xcolor}
\usepackage[colorlinks=true, urlcolor=blue!60!black, linkcolor=blue!60!black]{hyperref}
\setlength{\parindent}{2em}
\setlength{\parskip}{1.25em}
\renewcommand{\baselinestretch}{1.0}
\begin{document}
\begin{center}
{\LARGE \textit{Mahmoud v. Taylor}, 606 U.S. 522 (2025)\par}
\vspace{0.8em}
{\large Case Brief --- Four-Page Edition\par}
\end{center}
\vspace{1em}

\section{Memory Jogger}

Ruling 6--3 that parents were entitled to a preliminary injunction because they were likely to succeed on their Free Exercise challenge to mandatory use of LGBTQ+-inclusive storybooks without notice or opt-outs, the Court applied \textit{Wisconsin v.\ Yoder} to trigger strict scrutiny without asking whether the policy was neutral and generally applicable (\href{https://www.law.cornell.edu/supremecourt/text/24-297}{Cornell LII}).

\section{Detailed Case Facts}

In 2022--2023 the Montgomery County, Maryland, Board of Education, serving more than 160,000 students in a highly religiously diverse county, introduced five ``LGBTQ+-inclusive'' storybooks into its kindergarten-through-fifth-grade English curriculum, including \textit{Prince \& Knight}, \textit{Uncle Bobby's Wedding}, and \textit{Born Ready}. Teacher guidance encouraged educators to ``[d]isrupt the either/or thinking'' of students and to treat comments questioning same-sex marriage or gender identity as ``hurtful.'' The Board initially gave notice and religious opt-outs, but in March 2023 it rescinded them, citing disruption, absenteeism, and stigma, while still allowing opt-outs from its sex-education unit.

\section{Procedural History}

Muslim, Catholic, and Ukrainian Orthodox parents, including Tamer Mahmoud and Enas Barakat, and the group Kids First sought a preliminary injunction restoring notice and opt-outs; two families had withdrawn their sons to private school. The District of Maryland denied relief, and a divided Fourth Circuit affirmed, finding no showing of coercion on a ``threadbare'' record. The Supreme Court reversed on June 27, 2025.

\section{Judicial Votes}

Reversed and remanded, \textbf{6--3}.
\begin{itemize}
\item \textbf{Majority:} Justice Alito, joined by Chief Justice Roberts and Justices Thomas, Gorsuch, Kavanaugh, and Barrett.
\item \textbf{Concurring:} Justice Thomas.
\item \textbf{Dissenting:} Justice Sotomayor, joined by Justices Kagan and Jackson.
\end{itemize}

\section{Holding}

Denying parents notice and religious opt-outs from normative elementary instruction on gender and sexuality ``substantially interferes with the religious development'' of their children under \textit{Yoder}, triggers strict scrutiny regardless of neutrality or general applicability under \textit{Employment Division v.\ Smith}, and fails narrow tailoring. The parents are entitled to a preliminary injunction requiring advance notice and opt-outs.

\section{Analysis of Opinions}

\textbf{Majority (Alito).} Parents have a right ``to direct `the religious upbringing' of their children.'' Beyond \textit{Barnette}'s direct compulsion, \textit{Yoder} bars policies that create ``pressure to conform'' and pose ``a very real threat of undermining'' parental instruction. Given the children's ages (5 to 11), the books were ``unmistakably normative,'' presenting same-sex marriage as cause for celebration and endorsing fluid gender identity; paired with guidance to ``disrupt'' student thinking, they posed an ``objective danger.'' \textit{Bowen v.\ Roy} and \textit{Lyng} concern internal government affairs; \textit{Yoder} is not confined to the Amish; plaintiffs need not await harm; and private schooling cannot cure an unconstitutional condition on public education. Burdens ``of the same character as that in \textit{Yoder}'' go straight to strict scrutiny. The policy failed it: the Board already allowed other opt-outs, its burdens were ``self-inflicted,'' and ``[t]here is no \textit{de maximis} exception to the Free Exercise Clause.''

\textbf{Thomas (concurring).} Elementary gender-identity instruction lacks historical pedigree; like the conformity struck down in \textit{Pierce}, weaving themes into core classes would ``give schools a playbook for evading the First Amendment.''

\textbf{Sotomayor, with Kagan and Jackson (dissenting).} The Clause bars coercion, not mere ``exposure''; \textit{Lyng} and \textit{Bowen} foreclose the ``threat of undermining'' test; bypassing \textit{Smith} distorts precedent; and individualized opt-outs will cause administrative ``chaos'' and chill curricula.

\textbf{Later developments.} In \textit{Mirabelli v.\ Bonta} (2026), relying on \textit{Mahmoud}, the Court held parents likely to succeed against California policies barring staff from disclosing a student's social transition, a ``greater'' intrusion than the storybooks.

\section{Examples: Future Applications}

\begin{enumerate}
\item \textbf{Same-Side: Workbook Hostile to Religion.} A mandatory second-grade workbook, \textit{No Gods Needed}, has teachers call Scripture ``just myths,'' with no opt-out. It substantially interferes and fails tailoring.
\item \textbf{Same-Side: Lessons Against Religious Practices.} A kindergarten ``fairness'' unit tells children to rethink religious dietary and modesty rules. The burden is ``of the same character as that in \textit{Yoder}.''
\item \textbf{Opposite-Side: High-School Evolution.} A neutral eleventh-grade biology text teaches evolution. Older students, neutral content, and a compelling interest favor the district.
\item \textbf{Opposite-Side: Removing Books for All.} Parents sue to remove books for every student. \textit{Mahmoud} protects only their own children and gives no curricular veto.
\item \textbf{Fence-Sitter: Seventh-Grade Health Lesson.} A factual unit on gender identity requires respect but does not correct religious views. Twelve-year-olds are less impressionable, and the line between ``exposure'' and interference is unsettled.
\end{enumerate}

\section{Critique}

\textbf{Erosion of \textit{Smith} without overruling it.} By treating \textit{Yoder} as a standalone trigger for strict scrutiny without a ``hybrid'' due-process claim, \textit{Mahmoud} narrows \textit{Smith} substantially while leaving lower courts uncertain which other Free Exercise burdens qualify as ``\textit{Yoder}-like.''

\textbf{Line-drawing and administrability.} Distinguishing impermissibly ``normative'' instruction from neutral ``exposure'' forces judges to parse children's books and teacher scripts, raising Justice Sotomayor's concern that resource-strapped districts may self-censor curricula rather than manage individualized opt-outs.

\section{Key Quotations}

\begin{quote}
``When the burden imposed is of the same character as that imposed in \textit{Yoder}, we need not ask whether the law at issue is neutral or generally applicable before proceeding to strict scrutiny. That much is clear from our decisions in \textit{Yoder} and \textit{Smith}.''
\end{quote}
\noindent --- Justice Alito, Opinion of the Court (slip op., at 36).

\begin{quote}
``A government burdens the religious exercise of parents when it requires them to submit their children to instruction that poses `a very real threat of undermining' the religious beliefs and practices that the parents wish to instill. \textit{Wisconsin v. Yoder}, 406 U. S. 205, 218 (1972). And a government cannot condition the benefit of free public education on parents' acceptance of such instruction.''
\end{quote}
\noindent --- Justice Alito, Opinion of the Court (slip op., at 1--2).

\begin{quote}
``When it comes to instruction that would burden the religious exercise of parents, the Board cannot escape its obligations under the Free Exercise Clause by crafting a curriculum that is so burdensome that a substantial number of parents elect to opt out. There is no de maximis exception to the Free Exercise Clause.''
\end{quote}
\noindent --- Justice Alito, Opinion of the Court (slip op., at 39).

\end{document}
